\subsection{Different designation origins require their own gain histories}
\label{sec:origin-matched-cessation}

A reduction of five tonnes does not identify a reserve when the tonnes entered
at different prices. Consider two forecast purchase origins, A and B, each of
ten tonnes. A was designated at a price of 8 and B later at a price of 9.
Each has its own documented match to an external linear instrument from the
date of designation. The instruments entered at zero value and cost. Under
the supplied model there is no discounting, carry, credit or basis difference,
and each instrument's gain exactly offsets its named item's risk change.
Prices are currency units per tonne; gains and all balances below are currency
units. These are synthetic measurement assumptions, rather than a procedure
for deriving actual forward values from spot prices.

At a subsequent measurement price of 11, A's instrument has gained
\(10(11-8)=30\), while B's has gained \(10(11-9)=20\).
The relationship's reserve is 50. If five tonnes from A cease to qualify, their
recorded effective gain is \(5(11-8)=15\). If five from B cease to qualify,
it is \(5(11-9)=10\). In both cases the stopped quantity is one quarter of
the total. Multiplying 50 by that nominal fraction gives 12.5, which is wrong
for either identified history. Quantity without the designation origin has
lost information needed for this calculation.

The accounting basis remains the historical IFRS 9 text in EU 2023/1803
\citep{eu20231803}. Paragraph B6.5.17 preserves an added volume's own
designation date and reference price; paragraph B6.5.20 leaves the measurement
of continuing quantities unaffected by removal. Paragraphs 6.5.6 and
B6.5.22--28 make cessation prospective when qualification fails.
The example supplies the required documentation, individual eligibility,
management as a group and proportional variability under paragraphs 6.4.1
and 6.6.1--3. It also supplies the finding that the stopped forecast volume
is no longer highly probable and that the rest continues to qualify.
Choosing a quantity in a report does not establish those facts.

The retained gain histories support a narrower deduction than a general
allocation of a heterogeneous reserve. Let origin \(j\) have quantity
\(q_j\), designation date \(d_j\) and original price \(f_j\).
At a supplied measurement price \(S\), its cumulative instrument gain and
the opposite change in its hedged purchase cash flows are
\[
 G_j=q_j(S-f_j),\qquad H_j=-G_j.
\]
Only origins already designated at that date enter the sums. All active
origins' gains have the same weak sign at each supplied measurement, and
there is no historical ineffectiveness. For the relationship as a whole,
\(D=\sum_jG_j\) and \(H=\sum_jH_j=-D\). Consequently the cumulative
signed lower-of calculation required by paragraph 6.5.11 gives
\[
 R=\operatorname{sgn}(D)\min\{|D|,|H|\}=D.
\]
At zero, both reserve and ineffectiveness are zero. The independently
documented, fully effective origin gains therefore reconcile to the
relationship's reserve. This reasoning retains the relationship calculation;
it does not substitute a sum of independent lower-of measurements. The match
must have existed from designation, rather than be invented later to explain
an aggregate reserve. The derivation establishes one traceable attribution
for this exact history, without establishing uniqueness for general histories.

Suppose quantity \(z\), strictly between zero and \(q_k\), stops from named
origin \(k\) at a same-date measured price \(S_d\). Keep the actual
instruments. Denote the frozen gross gain on that portion by \(F_0\), and
the continuing quantities by \(q'_k=q_k-z\) and \(q'_j=q_j\) otherwise.
Their original dates and prices remain in force. At a later price \(S\),
\begin{align*}
 F_0&=z(S_d-f_k),\\
 R_c&=\sum_j q'_j(S-f_j),\\
 U&=z(S-S_d).
\end{align*}
Here \(R_c\) is the continuing relationship's reserve, again obtained from
its exact offset and cumulative lower-of test. The stopped instrument
portion's change \(U\) after cessation enters profit under paragraph 5.7.1;
it cannot enlarge the stopped reserve. The identity follows by splitting
the original \(k\)-origin quantity:
\[
 q_k(S-f_k)=(q_k-z)(S-f_k)+z(S_d-f_k)+z(S-S_d).
\]
Adding the other origins gives total actual derivative value
\(D=R_c+F_0+U\).

To follow later expectation or loss findings, let \(L\) be cumulative signed
reserve released to profit. A released gain makes \(L\) positive; a released
loss makes it negative. Remaining stopped reserve \(F\) and cumulative profit
\(P\) then satisfy
\[
 F=F_0-L,\qquad P=U+L,\qquad D=R_c+F+P.
\]
The distinction between no longer highly probable and no longer expected
remains the one in paragraph 6.5.12. Expected cash flows retain \(F\),
subject to the immediate unrecoverable-loss requirement in
paragraph 6.5.11(d)(iii). Cancellation releases the entire remaining \(F\)
once. If a loss was released earlier, only the balance remains available
for cancellation or an eventual inventory transfer.

Follow the two origins to a terminal price of 12. A's actual instrument
settles for 40 and B's for 30 in every history; neither instrument was sold
when hedge accounting stopped. If five tonnes from A were stopped and then
canceled, the continuing reserve is
\(5(12-8)+10(12-9)=50\). The stopped instrument's post-cessation gain is
\(5(12-11)=5\), and releasing its frozen gain of 15 gives profit 20.
If instead five tonnes from B were stopped and canceled, continuing reserve
is \(10(12-8)+5(12-9)=55\). Releasing its frozen gain of 10 gives profit 15.
For example, canceling A debits the stopped reserve by 15 and credits profit
by 15; the later gain of 5 is a separate instrument remeasurement.

With opening cash and capital of 300, both canceled histories buy fifteen
tonnes for 180. Paragraph 6.5.11(d)(i) transfers the attached reserve into
the nonfinancial asset's initial cost. Supplier cash and the two actual
instrument settlements remain gross:
\begin{center}
\small
\begin{tabular}{lrr}
\toprule
& Stop A, then cancel & Stop B, then cancel\\
\midrule
Continuing reserve before purchase &50&55\\
Profit &20&15\\
Supplier payment &180&180\\
Actual instrument receipts &70&70\\
Inventory cost &130&125\\
Final cash &190&190\\
Cash plus inventory &320&315\\
\bottomrule
\end{tabular}
\end{center}
The identical derivative cash does not make the histories interchangeable.
Their different original effective gains have different consequences when
the associated forecast purchases are canceled.

If the stopped purchases remain expected and occur instead, twenty tonnes
cost 240. Stopping A leaves reserve \(50+15=65\) attached; stopping B leaves
\(55+10=65\). Both then have inventory cost 175, cash
\(300-240+70=130\), and profit 5. The equality of these final totals is
consistent with different intermediate reserve origins. An implementation
must retain those origins because later cancellation or loss findings act
on the named stopped amount.

The same signs work for losses. With the original prices still 8 and 9,
stop five tonnes from A at price 7. Their frozen reserve is \(-5\).
A supplied finding that 2 will not be recovered releases \(-2\) to profit
and leaves \(-3\) attached. If the purchase remains expected and the
terminal price is 6, continuing reserve is
\(5(6-8)+10(6-9)=-40\), the stopped instrument's later gain is
\(5(6-7)=-5\), and profit is \(-5-2=-7\).
Purchasing twenty tonnes for 120 and attaching reserve \(-43\) gives
inventory 163. Actual instruments settle at \(-20-30=-50\), leaving
cash 130 and net assets 293, equal to capital less the loss of 7.
If cancellation follows the partial loss release, only the remaining loss
of 3 moves to profit. A subsequent rise in instrument value does not restore
the 2 already released from the stopped reserve.

More generally, let \(Q\) be continuing purchase quantity plus \(z\) only
if the stopped purchase still occurs. Immediately before purchase, the
attached reserve is \(B=R_c+F\). For opening cash and capital \(C_0\),
\[
 I=QS-B,\qquad C=C_0-QS+D,\qquad
 C+I=C_0+(D-B)=C_0+P.
\]
Thus the inventory basis \(I\), final cash \(C\) and profit reconcile
without using net settlement to conceal financing. The finite example pays
the supplier before settling each instrument in origin order, so every
intermediate cash balance must also be nonnegative. Capital of 240 exactly
funds the retained twenty-tonne purchase at 12; 239 does not, even though
the later instrument receipts would make final cash positive. This payment
order is a supplied convention, rather than a universal legal requirement.

The associated read-only method, \texttt{OriginMatchedCessation/1}, retains
the item and instrument identities, original dates and prices, each supplied
measurement, signed releases and the final basis transfer. It proposes
balanced journals without posting them. Missing qualification or matching
facts leave the report unresolved; contradictory values are rejected.
Historical ineffectiveness, opposed active-origin gains, acquired nonzero
origins, additional designations, repeated stops and a failure of continuing
loss recoverability lie outside this finite method. A negative resulting
inventory basis is also refused: the accounting for that extension has not
been established here. General heterogeneous attribution remains open.
Paragraph B6.6.14 addresses allocation across income-statement lines and does
not supply the missing stopped-reserve measurement rule.
